\documentclass[12pt]{ximera}
\title{Homework 9: Huntington-Hill, Paradoxes, and the Subtraction Game}
\begin{document}
\begin{abstract}
Sections 4.5--4.6 and the subtraction game: the Huntington-Hill method with geometric
mean cutoffs, the new-states paradox, and a subtraction game the second player wins.
\end{abstract}
\maketitle

\section*{Sections 4.5--4.6: Huntington-Hill and the Paradoxes}

Round all apportionment answers to three decimal places.

\begin{problem}
A city school district is apportioning $40$ teacher's aides to its five schools.

\begin{center}
\begin{tabular}{|r|c|c|c|c|c|c|}\hline
\textbf{School} & A & B & C & D & E & \textbf{Total} \\\hline
\textbf{Population} & 185 & 261 & 380 & 862 & 1312 & 3{,}000 \\\hline
\end{tabular}
\end{center}

Solve this apportionment using the Huntington-Hill method, rounding based on comparing
each quota to the geometric mean $\sqrt{\lfloor q_i\rfloor \cdot \lceil q_i\rceil}$ of
its lower and upper quotas.

\begin{prompt}
\textbf{(a)} Compute the standard divisor and the standard quotas.

Standard divisor: $\answer{75}$

\smallskip
$q_A = \answer[tolerance=0.0005]{2.467}$ \quad
$q_B = \answer[tolerance=0.0005]{3.48}$ \quad
$q_C = \answer[tolerance=0.0005]{5.067}$ \quad
$q_D = \answer[tolerance=0.0005]{11.493}$ \quad
$q_E = \answer[tolerance=0.0005]{17.493}$

\begin{feedback}
\[ SD = \frac{3000}{40} = 75, \]
so the quotas are
\[
q_A = 2.467,\quad q_B = 3.480,\quad q_C = 5.067,\quad q_D = 11.493,\quad q_E = 17.493,
\]
totaling $40$. \checkmark
\end{feedback}
\end{prompt}

\begin{prompt}
\textbf{(b)} Compute each geometric mean cutoff, and round accordingly.

Cutoffs: $A: \answer[tolerance=0.0005]{2.449}$ \quad
$B: \answer[tolerance=0.0005]{3.464}$ \quad
$C: \answer[tolerance=0.0005]{5.477}$ \quad
$D: \answer[tolerance=0.0005]{11.489}$ \quad
$E: \answer[tolerance=0.0005]{17.493}$

\smallskip
Rounded quotas: $A: \answer{3}$ \quad $B: \answer{4}$ \quad $C: \answer{5}$ \quad
$D: \answer{12}$ \quad $E: \answer{18}$ \quad giving a total of $\answer{42}$

\begin{hint}
Round \emph{up} when the quota exceeds its geometric mean cutoff, and \emph{down}
otherwise. Several of these are very close calls --- keep all three decimals.
\end{hint}

\begin{feedback}
\[
\begin{array}{lllll}
A: & q = 2.467 & \sqrt{2\cdot 3}=2.449 & 2.467 > 2.449 & \to 3\\
B: & q = 3.480 & \sqrt{3\cdot 4}=3.464 & 3.480 > 3.464 & \to 4\\
C: & q = 5.067 & \sqrt{5\cdot 6}=5.477 & 5.067 < 5.477 & \to 5\\
D: & q = 11.493 & \sqrt{11\cdot 12}=11.489 & 11.493 > 11.489 & \to 12\\
E: & q = 17.493 & \sqrt{17\cdot 18}=17.493 & \text{just above} & \to 18
\end{array}
\]
D and E are genuinely close: $q_D = 11.4933$ against a cutoff of $11.4891$, and
$q_E = 17.49333$ against $17.49286$. Both round up, but only barely.

The total is $3+4+5+12+18 = 42$, which is $2$ too many. A modified divisor is needed.
\end{feedback}
\end{prompt}

\begin{prompt}
\textbf{(c)} Which modified divisor should be used?

\begin{multipleChoice}
\choice{$74.75$}
\choice[correct]{$75.25$}
\end{multipleChoice}

With it, the final apportionment is \\
$A: \answer{3}$ \quad $B: \answer{4}$ \quad $C: \answer{5}$ \quad $D: \answer{11}$
\quad $E: \answer{17}$

\begin{hint}
The standard divisor gave too many seats. Does a larger or a smaller divisor reduce
the quotas?
\end{hint}

\begin{feedback}
Since the standard divisor gave $42$ seats --- two too many --- the quotas must come
down, which means a \emph{larger} divisor. So $75.25$ is the right choice; $74.75$
would only make things worse (it still totals $42$).

With $d = 75.25$:
\[
\begin{array}{lllll}
A: & q = 2.458 & \text{cutoff } 2.449 & \text{above} & \to 3\\
B: & q = 3.468 & \text{cutoff } 3.464 & \text{above} & \to 4\\
C: & q = 5.050 & \text{cutoff } 5.477 & \text{below} & \to 5\\
D: & q = 11.455 & \text{cutoff } 11.489 & \text{below} & \to 11\\
E: & q = 17.435 & \text{cutoff } 17.493 & \text{below} & \to 17
\end{array}
\]
which totals $40$. \checkmark

The two borderline cases, D and E, are precisely the ones that flip.

\textbf{Compare with Webster's method}, which uses conventional rounding. At the
standard divisor of $75$ that would give $2, 3, 5, 11, 17 = 38$, so a smaller divisor
would be needed --- and school A, whose quota is $2.467$, would be rounded \emph{down}
to $2$ rather than up to $3$. Huntington-Hill's geometric-mean cutoff sits slightly
below the arithmetic midpoint, which nudges small districts upward. That is exactly why
it is the method the U.S.\ House of Representatives uses.
\end{feedback}
\end{prompt}
\end{problem}

\begin{problem}
Narnia has grown from two states to three. To accommodate the new state of
Chippingford, the Narnian congress has been increased by $5$ seats and the new total of
$95$ seats has been reapportioned:

\begin{center}
\begin{tabular}{|r|c|c|c|}\hline
\textbf{State} & Beruna & Beaversdam & Chippingford \\\hline
\textbf{Population} & 7600 & 9720 & 1000 \\\hline
\textbf{Original apportionment} & 39 & 51 & -- \\\hline
\textbf{New apportionment} & 40 & 50 & 5 \\\hline
\end{tabular}
\end{center}

\begin{prompt}
\textbf{(a)} Why might Beaversdam residents be frustrated with this reapportionment?

\begin{freeResponse}
\end{freeResponse}

\begin{feedback}
Beaversdam \emph{lost} a seat, dropping from $51$ to $50$, even though:
\begin{itemize}
\item its own population did not change;
\item the congress got \emph{bigger}, gaining $5$ seats; and
\item the new seats were added specifically to accommodate Chippingford, which took
      exactly the $5$ seats that were added.
\end{itemize}
Worse, the seat Beaversdam lost went to Beruna --- a state whose population also did
not change. Adding a new state with its own new seats should not shuffle
representation between two untouched states.
\end{feedback}
\end{prompt}

\begin{prompt}
\textbf{(b)} What is the formal term for this phenomenon?

\begin{multipleChoice}
\choice{The Alabama paradox}
\choice{The population paradox}
\choice[correct]{The new-states paradox}
\choice{A violation of the quota rule}
\end{multipleChoice}

\begin{feedback}
The \textbf{new-states paradox}: adding a new state, along with the seats to
accommodate it, changes the apportionment among the existing states.

The neighboring paradoxes are worth distinguishing. The \emph{Alabama paradox} is when
increasing the total number of seats causes a state to lose one. The \emph{population
paradox} is when a state whose population is growing loses a seat to one that is
shrinking. All three afflict the same method.
\end{feedback}
\end{prompt}

\begin{prompt}
\textbf{(c)} Of the apportionment methods we've discussed, which must the Narnian
congress be using?

\begin{multipleChoice}
\choice[correct]{Hamilton's method}
\choice{Jefferson's method}
\choice{Adams' method}
\choice{Webster's method}
\end{multipleChoice}

\begin{feedback}
\textbf{Hamilton's method}. All three apportionment paradoxes --- Alabama, population,
and new-states --- are possible only under Hamilton's method.

The divisor methods (Jefferson, Adams, Webster, Huntington-Hill) are immune to all
three paradoxes, because they apportion by picking a divisor and rounding, with no
step that reshuffles surplus seats based on fractional parts. The trade-off is that
divisor methods can violate the quota rule, which Hamilton's method never does. As the
Balinski-Young theorem proves, no method can avoid both problems at once.
\end{feedback}
\end{prompt}
\end{problem}

\begin{problem}
Alice and Bob are playing a version of the Subtraction Game where they start with a pile
of $54$ stones and on each turn they may take $1$, $2$, $3$, $4$, or $5$ stones. The
objective is to claim the last stone.

\begin{prompt}
\textbf{(a)} Describe the winning strategy.

\begin{freeResponse}
\end{freeResponse}

\begin{feedback}
Always leave your opponent a multiple of $6$: whenever they take $k$ stones, take $6-k$.
Each round removes exactly $6$, so your opponent eventually faces $0$.
\end{feedback}
\end{prompt}

\begin{prompt}
\textbf{(b)} If both players play optimally and Alice goes first, who will win?

\begin{multipleChoice}
\choice{Alice}
\choice[correct]{Bob}
\end{multipleChoice}

\begin{feedback}
Bob. $54=6\cdot 9$ is already a multiple of $6$, so Alice starts in a losing position:
whatever she takes, Bob restores a multiple of $6$.
\end{feedback}
\end{prompt}
\end{problem}

\end{document}
